% Generated draft: requires mathematical and source review.
\documentclass[11pt]{article}
\usepackage[margin=25mm]{geometry}
\usepackage{fontspec}
\setmainfont{Latin Modern Roman}
\usepackage{amsmath,amssymb,mathtools,unicode-math}
\setmathfont{Latin Modern Math}
\usepackage{enumitem,needspace,xurl,hyperref}
\setlength{\parindent}{0pt}
\setlength{\parskip}{5pt}
\setlength{\emergencystretch}{3em}
\newcommand{\N}{\mathbb{N}}
\newcommand{\Z}{\mathbb{Z}}
\newcommand{\Q}{\mathbb{Q}}
\newcommand{\R}{\mathbb{R}}
\newcommand{\C}{\mathbb{C}}
\newcommand{\F}{\mathbb{F}}
\newcommand{\A}{\mathbb{A}}
\newcommand{\PP}{\mathbb{P}}
\newcommand{\E}{\mathbb{E}}
\newcommand{\eps}{\varepsilon}
\newcommand{\dd}{\,\mathrm{d}}
\begin{document}
\section*{UnsolvedMath ID: 800272}
\textbf{Source problem code:} catalogue entry
\subsection{Definitions and mathematical statement}
Let $\PP^k(\C)$ denote complex projective space of complex dimension $k$. Consider a holomorphic self-map $f\colon\PP^k(\C)\to\PP^k(\C)$ and a fixed point $p$, so that $f(p)=p$. The source also asks for the local version, in which $f$ is defined only on an open neighborhood $U$ of $p$ in $\PP^k(\C)$ and satisfies $f(p)=p$. In local coordinates centered at $p$, this gives a holomorphic germ $F\colon(\C^k,0)\to(\C^k,0)$ with $F(0)=0$. Two representatives define the same germ if they agree on some neighborhood of $0$. Iteration is considered wherever the successive compositions are defined; no invariant neighborhood is assumed.

\textbf{Classification problem.} Classify the dynamics near the fixed point $p$ of such a holomorphic map, both for maps defined on all of $\PP^k(\C)$ and for maps defined only near $p$.

\textbf{Unresolved scope.} The source does not specify a range of dimensions $k$, an equivalence relation for the classification, or the dynamical features to be classified. In particular, it does not say whether classification is intended up to holomorphic conjugacy, topological conjugacy, or another relation. It imposes no stated hypotheses on the derivative at the fixed point or on local invertibility. The local formulation above interprets the neighborhood version as a holomorphic self-map germ of the same ambient complex manifold; the local codomain is not separately stated in the source. Thus the supplied request is a broad classification program rather than a uniquely specified theorem or conjecture.
\subsection{Short English statement}
Describe and classify the behavior near a fixed point of a holomorphic map of complex projective space, including the case where the map is defined only near that point. The source leaves the meaning and scope of classification unspecified. Its supplied status is \emph{open}; that label has not been independently verified.
\subsection{Sources}
\begin{itemize}
\item Local UnsolvedMath registry, permanent ID \textbf{800272}, exact source problem code \texttt{catalogue entry}, title \emph{Several Complex Variables Problem List --- Problem 71}. The record supplies the status \emph{open} and lists Buzzard, Fornaess, and Krantz under proposer attribution, with year 1995.
\item Buzzard, Fornaess, and Krantz, \emph{The Several Complex Variables Problem List} (1995), Problem 71. Supplied source URL: \url{https://arxiv.org/abs/math/9509202}. The record reports extraction by \texttt{arxiv\_latex} and an access date of \texttt{2026-10-16}; these are supplied provenance details, not independently verified access or a current literature assessment.
\item The supplied problem text ends with the contact signature \texttt{sibony@anh.matups.fr}. This signature is retained as source-specific attribution without inferring a fuller author identity.
\end{itemize}
\end{document}
